% Fragmento integrable; no documento autónomo y no orden de compilación.
% Fecha: 2026-09-30. No modifica fuentes publicadas.
% Procedencia: reunión de RETROACCION_PCH_Y_RESTRICCIONES_VARIACIONALES.md,
% PRECUANTIZACION_PCH_Y_FIDELIDAD.md y CIERRE_CANONICO_TOTAL_Y_REPRESENTACION.md.
% Prerrequisitos de inserción: núcleo común APP--TRIT--TPK; VIII/30d,
% VIII/31 y VIII/33; lectores de acción, gamma y acoplamientos.
% Los nuevos labels usan exclusivamente el prefijo hmtcuatro:.
\section{Joint action, backreaction and canonical closure of the constraints}
\label{hmtcuatro:accion-retroaccion-restricciones}

The causal sequence of this development is
\[
 \mathrm{APP}\longrightarrow\mathrm{TRIT}\longrightarrow\mathrm{TPK}
 \longrightarrow\text{enriched state}
 \longrightarrow\text{joint discrete structure of the continuum}.
\]
The sheets, orientation, residues, carry, routes, boundary and memory of this construction are preserved. The action scale and the couplings are HMT outputs received before using variational and symplectic language. They are not selected by target values or by the curvature conclusion to be proved. Nor does phase return reset the enriched state.

The previously constructed geometric realization provides a nondegenerate cotetrad \(e\), its metric connection and the spinorial lift of transport. Here that geometry is varied together with the matter fields. We shall bring together in a single proof the common action, its torsional reduction, the Legendre transformation, the first-class property, constraint propagation and the characteristic flatness of the canonical connection resulting from the same action. This last realization is not identified by name with another interacting Hamiltonian.

\subsection{Starting structure and action with the full source}
\label{hmtcuatro:datos-accion-total}

Let \(M\) be a smooth oriented four-dimensional region with spin structure, signature \((-+++)\) and a regular spatial foliation. The cotetrad takes values in the length line of the received realization. We write
\[
 \Sigma^{IJ}=e^I\wedge e^J,\qquad
 F_\omega=\mathrm d\omega+\omega\wedge\omega,\qquad
 \mathsf P_\gamma=\star+\gamma^{-1}I,\qquad \star^2=-I.
\]
The star in this formula acts on internal bivectors; it is not exterior duality of forms. During variation, the received sections \(c>0,\hbar>0,G>0,\gamma\in\mathbb R\setminus\{0\}\) and \(\Lambda\) remain fixed. With \(\kappa=8\pi G/c^4\), the total action is
\begin{equation}
 \begin{split}
 S_{\rm tot}={}&\frac1{2\kappa c}
   \int_M\Sigma_{IJ}\wedge(\mathsf P_\gamma F_\omega)^{IJ}
 -\frac{\Lambda}{\kappa c}\int_M\operatorname{vol}_e\\
 &+S_{\rm YM}[e,A]+S_H[e,A,H]
   +S_W[e,\omega,A,\Psi]+S_Y[e,H,\Psi]+S_{\partial}.
 \end{split}
 \label{hmtcuatro:accion-total}
\end{equation}

\begin{definition}[Variational and canonical unification]
\label{hmtcuatro:def-unificacion}
HMT variational and canonical unification is the composition of the geometric, chromatic and electroweak realizations of the same enriched state in the common action \eqref{hmtcuatro:accion-total}, whose variations produce joint sources and whose reduction preserves a common constraint algebra and its propagation. The electromagnetic component belongs to the electroweak realization. Theorem~\ref{hmtcuatro:teorema-canonico-total} establishes this composition in the regular domain and under the stated boundary conditions.
\end{definition}
The electroweak couplings $g,g'$ and the self-coupling $\lambda_H$ are inherited from the construction in \ref{amp-higgs-seccion} and Theorem~\ref{amp-higgs-inversa-teorema}; the present variation holds these sections fixed. Comparison between sections is a distinct operation, preserved in the compatibility-horizon development.


The symbol \(H\) without a subscript denotes the Higgs field here, not a Hamiltonian. All spacetime contractions of matter use the same \(e\). The boundary problem is that received from the Cartan--Holst variation: interior variations, compatible fixed boundary data, or a completion cancelling the same boundary term.

To fix the content of \eqref{hmtcuatro:accion-total}, the internal connections act on the chiral bundle
\[
 (S_L\otimes E_L)\oplus(S_R\otimes E_R).
\]
That bundle is not replaced by a vector coupling. In the integer normalization \(y=6Y\), the received multiplets have weights \(1,4,-2,-3,-6\), respectively for \(q_L,u_R,d_R,\ell_L,e_R\); the doublet \(H\) has weight \(3\). Colour acts on the factor \(\mathbb C^3\) of quarks and trivially on leptons. The family factor and its maps \(Y_u,Y_d,Y_e\) are preserved. This block does not add a right-handed partner to a multiplet that has none.

With \(\widetilde H=i\sigma_2\overline H\), the material map is
\begin{equation}
 \begin{split}
 \mathsf Y_q(H)(u,d)&=
       \widetilde H\otimes Y_u u+H\otimes Y_d d,\\
 \mathsf Y_\ell(H)e&=H\otimes Y_e e .
 \end{split}
 \label{hmtcuatro:yukawa}
\end{equation}
The identity on the colour factor is understood in the first line. The action term is \(-(\bar\Psi_L\mathsf Y(H)\Psi_R+\mathrm{h.c.})\), with the received dimensional sections. Equivariance is not merely nominal: for \(U\in SU(2)\), \(i\sigma_2\overline U=U i\sigma_2\), so that \(\widetilde H\mapsto U\widetilde H\). The weights of the three columns are \(-3+4=1\), \(3-2=1\) and \(3-6=-3\). Colour is intertwined by the identity, and the family maps commute with the internal action. Therefore, for each internal transformation \(u\),
\begin{equation}
 \mathsf Y(uH)\rho_R(u)=\rho_L(u)\mathsf Y(H).
 \label{hmtcuatro:equivariancia-yukawa}
\end{equation}
If the preceding material realization also contains other multiplets, their maps are incorporated with their own equivariance identity; this proof invents no new states.

In a units chart \(\hbar=c=1\), the material density summarizing these actions has the form
\begin{equation}
 \begin{split}
 \mathcal L_m={}&-\frac14 F_{B,\mu\nu}F_B^{\mu\nu}
 -\frac14\sum_a F_{W,\mu\nu}^aF_W^{a,\mu\nu}
 -\frac14\sum_a F_{C,\mu\nu}^aF_C^{a,\mu\nu}\\
 &-(\nabla_\mu H)^\dagger\nabla^\mu H-V(H)
 +\mathcal L_W(\Psi_L,\nabla_L)+\mathcal L_W(\Psi_R,\nabla_R)\\
 &-\bigl(\bar\Psi_L\mathsf Y(H)\Psi_R+\mathrm{h.c.}\bigr).
 \end{split}
 \label{hmtcuatro:densidad-material}
\end{equation}
The curvatures include their couplings; these and the component normalizations are those of their source documents. Writing in that chart neither removes the action section nor recalibrates the coefficients. In particular, no new numerical value for the strong coupling is introduced here. The spin connection of \(\nabla_L,\nabla_R\) is that of \(e,\omega\); the connections \(B,W,C\) act in the representations just described. The constant of \(V\) is retained, since its volume variation belongs to the gravitational source.

The spinorial part preserves the explicit conventions of the preceding development: \(\{g^I,g^J\}=2\eta^{IJ}I\), \(\mathbb B_{IJ}=\frac14[g_I,g_J]\), \(A_D=-ig^0\), \(\bar\psi=\psi^\dagger A_D\) and \(D\psi=\mathrm d\psi+\frac12\omega^{IJ}\mathbb B_{IJ}\psi+A\psi\). Its symmetric kinetic term is
\[
 \frac{\hbar}{2}
 [\bar\psi g^I D\psi-(D\bar\psi)g^I\psi]\wedge\eta_I,
 \qquad \eta_I=\iota_{E_I}\operatorname{vol}_e .
\]
The sum runs over the fields actually present, with their chiral projections, not over a vector duplication of the catalogue. The Yukawa terms are algebraic and do not depend on \(\omega\). The total current is fixed by
\begin{equation}
 \delta_\omega S_m=-\frac12\int_M\delta\omega^{IJ}\wedge\sigma_{IJ},
 \qquad
 \sigma_{JK}=\sum_s-\frac{\hbar}{2}
   \bar\psi_s\{g^I,\mathbb B_{JK}\}\psi_s\,\eta_I.
 \label{hmtcuatro:corriente-total}
\end{equation}

\subsection{A proof of compatibility, propagation and closure}
\label{hmtcuatro:demostracion-unica}

\begin{theorem}[Canonical compatibility of the joint action]
\label{hmtcuatro:teorema-canonico-total} Consider \eqref{hmtcuatro:accion-total} on the preceding nondegenerate domain, with matter affine in the contorsion, and in a regular chart of the restricted Legendre transformation. Retain the first-order fermionic constraints and eliminate the second-class constraints through their Dirac bracket. For parameters with interior support, or with the compatible boundary without improper charges, the following hold jointly: the metric equation has the total material source of that same action; its diffeomorphism and gauge generators are first class; the constraints propagate on every regular coupled solution; and the total canonical potential induces a prequantum connection flat in the characteristic directions of the regular constraint surface.

The assertion concerns this action and this canonical realization. It does not automatically identify its connection with a different interacting Hamiltonian or prove its spatial--chiral limit.
\end{theorem}

\begin{proof}
\par\medskip\noindent\textbf{1. Torsional elimination without duplication.}\quad
Write \(\omega=\mathring\omega(e)+\mathfrak k\). The expansion \(F_\omega=F_{\mathring\omega}
+D_{\mathring\omega}\mathfrak k+\mathfrak k\wedge\mathfrak k\) and \(D_{\mathring\omega}\Sigma=0\) separate the boundary contribution
\[
 S_{\partial,\mathfrak k}
 =\frac1{2\kappa c}\int_{\partial M}
                  \Sigma_{IJ}\wedge(\mathsf P_\gamma\mathfrak k)^{IJ}
\]
from the algebraic density
\begin{equation}
 \mathcal Q_{e,\gamma}(\mathfrak k)
       -\frac12\mathfrak k^{IJ}\wedge\sigma_{IJ},\qquad
 \mathcal Q_{e,\gamma}(\mathfrak k)=\frac1{2\kappa c}
       \Sigma_{IJ}\wedge
       \mathsf P_\gamma(\mathfrak k\wedge\mathfrak k)^{IJ}.
 \label{hmtcuatro:cuadratica-contorsion}
\end{equation}
The original variation gives \(D_\omega(\mathsf P_\gamma\Sigma)=\kappa c\,\sigma\). The Holst inverse and the torsion--contorsion inverse from the preceding development are invertible for the declared coframe and \(\gamma\). They thus determine a unique \(\mathfrak k_*[e,\sigma]\), linear in \(\sigma\). The radial variation of \eqref{hmtcuatro:cuadratica-contorsion} at its stationary point gives \(2\mathcal Q(\mathfrak k_*)=\frac12\mathfrak k_*^{IJ}\wedge\sigma_{IJ}\). Consequently,
\begin{equation}
 S_{\rm red}=\frac1{2\kappa c}\int_M(R-2\Lambda)\operatorname{vol}_e
       +S_m^{\rm LC}
       -\frac14\int_M\mathfrak k_*^{IJ}\wedge\sigma_{IJ}
       +S_{\partial,\rm red}.
 \label{hmtcuatro:accion-reducida}
\end{equation}
The first Bianchi identity cancels the Levi--Civita Holst term in this expression. The chain rule preserves the equations of \(e,A,H,\Psi\), since the coefficient of \(\delta\mathfrak k_*\) vanishes by stationarity. Since \(\sigma=\sum_s\sigma_s\), the effective contribution includes \(\mathfrak k_*[\sigma_s]\wedge\sigma_t\) for \(s\ne t\). It is not replaced by squares of isolated species, nor is the same \(\mathfrak k\) simultaneously retained as an independent variable.

\par\medskip\noindent\textbf{2. Full source and backreaction.}\quad
Define the action tensor by varying all the material terms of \eqref{hmtcuatro:accion-reducida}:
\begin{equation}
 \delta_g S_{m,\rm red}=-\frac12\int_M
       \mathcal T^S_{\mu\nu}\,\delta g^{\mu\nu}\operatorname{vol}_e,
 \qquad
 E_{\mu\nu}:=G_{\mu\nu}+\Lambda g_{\mu\nu}
                   -\kappa c\,\mathcal T^S_{\mu\nu}.
 \label{hmtcuatro:fuente-completa}
\end{equation}
The metric equation is \(E_{\mu\nu}=0\). The factor \(\kappa c\) corresponds to the length cotetrad; if \(T^{\rm phys}=c\mathcal T^S\), we write \(\kappa T^{\rm phys}\). The matter fields and internal connections are also varied. Thus the fields determine the source, and the geometry solving the equation in turn determines their connections, volume and evolution. This is backreaction, not evolution on a fixed background. The variation includes the volume of the Higgs potential and the metric dependence of the torsional current, as well as its cross-terms.

\par\medskip\noindent\textbf{3. Restricted Legendre transform and canonical reduction.}\quad
The total canonical potential is obtained from \(\delta L=\mathcal E\delta z+\mathrm d\theta(\delta z)\) by integrating \(\theta\) over the spatial leaf and performing the second-class reduction. Its geometric part before that reduction is
\[
 \theta_{\rm geom}(\delta)=
   \frac1{2\kappa c}(\mathsf P_\gamma\Sigma)_{IJ}
                       \wedge\delta\omega^{IJ}.
\]
The decomposition \(F_{0a}=\dot\omega_a-D_a\omega_0\), together with the same decomposition of the material actions, produces
\begin{equation}
 S_{\rm can}=\int d\tau\,
 \bigl[\Theta_{\rm red}(\dot z)-\mathcal C[N]-\mathcal D[\xi]
       -\mathcal G_{\rm L}[\lambda]-\mathcal G_{\rm int}[\chi]\bigr]
       +S_{\partial,\rm red}.
 \label{hmtcuatro:legendre-total}
\end{equation}
The absent velocities of lapse, shift or temporal connections are not artificially inverted.

The algebraic contorsion has constraints \(\pi_{\mathfrak k}=0\) and \(\partial\mathcal Q/\partial\mathfrak k-\sigma/2=0\). Their bracket matrix has the block
\[
 \begin{pmatrix}0&-M\\ M^T&B\end{pmatrix},
 \qquad M=\frac{\partial^2\mathcal Q}{\partial\mathfrak k^2}.
\]
The torsional inverse makes \(M\) invertible, and hence this matrix. These are second-class constraints. We use
\begin{equation}
 \{F,G\}_D=\{F,G\}
       -\{F,\chi_a\}(C^{-1})^{ab}\{\chi_b,G\}.
 \label{hmtcuatro:corchete-dirac}
\end{equation}
For the variables that do not depend on the contorsion--momentum pair, the lower-right block of \(C^{-1}\) is zero. This elimination produces no additional bracket between them. The first-order spinorial constraints are retained separately; in the graded phase space, the graded bracket and even generators are used. \(\Theta_{\rm red}\) denotes the full result, not only its geometric part.

Legendre inversion in regular directions recovers the same action. For example, \(N[p^2/(2w)+wV+H_{\rm grad}]\), with \(w=\sqrt{\det q}\), gives \(p=(w/N)D_\tau\phi\), and by substitution
\[
 L_H=\frac{w}{2N}|D_\tau\phi|^2-NwV-NH_{\rm grad}.
\]
\(D_\tau\) preserves the shift and temporal connection; \(w\) is not frozen. For \(K_{ij}=(\dot q_{ij}-\mathcal L_\xi q_{ij})/(2N)\), the gravitational contribution to momentum and its inverse are
\begin{equation}
 \Pi^{ij}=\frac{\sqrt q}{2\kappa c}(K^{ij}-q^{ij}K),\qquad
 K_{ij}=\frac{2\kappa c}{\sqrt q}
                  (\Pi_{ij}-\tfrac12q_{ij}\Pi).
 \label{hmtcuatro:legendre-geometrica}
\end{equation}
They yield
\[
 \mathcal C_{\rm grav}
 =\frac{2\kappa c}{\sqrt q}
       (\Pi^{ij}\Pi_{ij}-\tfrac12\Pi^2)
  -\frac{\sqrt q}{2\kappa c}(R^{(3)}-2\Lambda).
\]
If the spinor chart contributes a material shift of the total momentum \(p\), we use \(\Pi=p-p_m\); \(p_m\) is not discarded. The indefinite character of the DeWitt form does not prevent this algebraic inverse. No velocity inverse is required of the first-order fermionic kinetic term.

\par\medskip\noindent\textbf{4. Noether--Legendre identity and first-class property.}\quad
The coincidence of the generators with the constraints of the same action can be seen before imposing its equations. Set \(a=(2\kappa c)^{-1}\), \(B=\mathsf P_\gamma\Sigma\), \(u^I=\iota_\zeta e^I\) and \(\lambda^{IJ}=\iota_\zeta\omega^{IJ}\). From \(\mathcal L_\zeta\omega=\iota_\zeta F+D_\omega\lambda\) we obtain
\begin{equation}
 \begin{split}
 j_\zeta^{\rm geom}
 ={}&\mathrm d(a B_{IJ}\lambda^{IJ})
   -a(D_\omega B_{IJ})\lambda^{IJ}\\
 &-2a u^Ie^J\wedge(\mathsf P_\gamma F)_{IJ}
       +2a\Lambda u^I\eta_I .
 \end{split}
 \label{hmtcuatro:noether-legendre}
\end{equation}
The material part adds exactly its contributions to the equations of \(e,\omega,A\), its vertical generators and its boundary. Decomposing the same expression on a spatial leaf gives \eqref{hmtcuatro:legendre-total}. Stationary reduction and Dirac reduction transport these transformations instead of replacing them by momenta of an external action.

With the usual canonical convention for the constraints, the total brackets are, modulo the preserved Gauss generators,
\begin{equation}
 \begin{aligned}
 \{\mathcal D[\xi],\mathcal D[\eta]\}_D
    &\simeq\mathcal D[[\xi,\eta]],\\
 \{\mathcal D[\xi],\mathcal C[N]\}_D
    &\simeq\mathcal C[\mathcal L_\xi N],\\
 \{\mathcal C[N],\mathcal C[M]\}_D
    &\simeq\mathcal D[
       q^{ij}(N\partial_jM-M\partial_jN)\partial_i].
 \end{aligned}
 \label{hmtcuatro:algebra-restricciones}
\end{equation}
Indeed, tangential transformations have the Lie bracket and transport the lapse as a scalar. For normal transformations, differentiating their orthogonality to the tangent vectors and \(g(n,n)=-1\) determines the tangential variation of \(n\); with \(K_{ij}\) from \eqref{hmtcuatro:legendre-geometrica} we obtain \(\delta_N n=\operatorname{grad}_qN\). The deformation commutator, converted to the canonical bracket with these conventions, gives the last line of \eqref{hmtcuatro:algebra-restricciones}. Covariantized transport additionally preserves frame and gauge rotations. The transformations are tangent to the constraint surface of the invariant action, yielding the first-class property. The functions \(q^{ij}\) depend on the state and are not frozen. With interior support, no central boundary charge is added; an improper generator with nonzero charge is not declared a zero constraint.

\par\medskip\noindent\textbf{5. Propagation in the geometry that responds to matter.}\quad
Diffeomorphism invariance of \(S_{m,\rm red}\), evaluated on its matter equations, gives by integration by parts \(\mathring\nabla^\mu\mathcal T^S_{\mu\nu}=0\). The identity uses the full tensor of \eqref{hmtcuatro:fuente-completa}; it does not postulate separate conservation of each contribution. Bianchi then implies \(\mathring\nabla_\mu E^{\mu\nu}=0\), even before imposing all the metric equations.

In a local Gaussian chart, \(ds^2=-d\tau^2+q_{ij}dx^idx^j\), impose the evolution equations \(E_{ij}=0\) and set \(C=E_{00}\), \(C_i=E_{0i}\), \(K_{ij}=\dot q_{ij}/2\), \(K=q^{ij}K_{ij}\). Here \(K\) is only the trace of the second fundamental form. The two components of the divergence identity are
\begin{equation}
 \partial_\tau C-D_iC^i+KC=0,\qquad
 \partial_\tau C_i+KC_i=0.
 \label{hmtcuatro:propagacion}
\end{equation}
To check the signs, we use \(E^{00}=C\), \(E^{0i}=-C^i\), \(E^{ij}=0\), \(\Gamma^i{}_{0j}=K^i{}_j\) and \(\Gamma^\mu{}_{\mu0}=K\). On lowering the index of the spatial equation, \(\dot q_{ij}=2K_{ij}\) cancels the two connection terms. If initially \(C_i=0\), its homogeneous equation maintains \(C_i=0\); the first becomes \(\dot C+KC=0\) and preserves \(C=0\). \(q\) and \(K\) are those of the coupled solution. The Gaussian choice proves the local assertion; its tensorial form allows lapse and shift to be changed. Likewise, the Noether gauge identity relates the covariant divergence of the gauge equation to the matter equations and propagates Gauss when the spatial equations hold. Neither global existence nor absence of singularities for every solution is inferred.

\par\medskip\noindent\textbf{6. The same action induces characteristic closure.}\quad
Let \(\mathcal P_{\rm red}\) be the regular canonical chart obtained and \(\mathcal Z=\{C_A=0\}\) its regular surface of total constraints, including Gauss. We use a symbol distinct from \(\Sigma^{IJ}\) for this surface. The potential \(\Theta=\Theta_{\rm red}\) of \eqref{hmtcuatro:legendre-total} is retained, including the material and mixed terms. We fix
\begin{equation}
 \Omega=-\mathrm d_{\mathcal P}\Theta,\qquad
 \iota_{X_f}\Omega=\mathrm d_{\mathcal P}f,\qquad
 \{f,g\}_D=\Omega(X_f,X_g).
 \label{hmtcuatro:convenciones-simpecticas}
\end{equation}
In a canonical pair, \(\Theta=p\,\mathrm dq\) gives \(X_f=f_p\partial_q-f_q\partial_p\), \(\{q,p\}_D=1\), \(X_f(g)=-\{f,g\}_D\) and \([X_f,X_g]=-X_{\{f,g\}_D}\). In the graded phase space, this expression applies to the relevant even functions and constraints.

On the phase line trivialized over this chart, we construct
\begin{equation}
 \nabla_X=X-\frac{i}{\hbar}\Theta(X)
       =\delta_X+\frac{i}{\hbar}H_X^{\rm can},
 \qquad H_X^{\rm can}=-\Theta(X).
 \label{hmtcuatro:conexion-canonica}
\end{equation}
No arbitrary conjugation is chosen to make it flat: its coefficient comes from the action potential. For any fields \(X,Y\) of the chart, calculation gives
\begin{equation}
 \begin{aligned}
 \mathcal F^{\rm can}_{X,Y}
 &=XH_Y^{\rm can}-YH_X^{\rm can}-H_{[X,Y]}^{\rm can}
       +\frac{i}{\hbar}[H_X^{\rm can},H_Y^{\rm can}]\\
 &=-X\Theta(Y)+Y\Theta(X)+\Theta([X,Y])\\
 &=-\mathrm d_{\mathcal P}\Theta(X,Y)=\Omega(X,Y).
 \end{aligned}
 \label{hmtcuatro:curvatura-calculada}
\end{equation}
The \(H_X^{\rm can}\) are multipliers on the line, so their commutator is zero; derivatives of the coefficients and fields are not suppressed.

The already proved first-class property means \(\{C_A,C_B\}_D=f_{AB}{}^D C_D\). For \(V\) tangent to \(\mathcal Z\), \(\Omega(X_{C_A},V)=\mathrm d_{\mathcal P}C_A(V)=0\). The constraint fields are therefore characteristic. Under the declared regularity they generate the characteristic distribution of the coisotropic surface. Tensoriality of \(\Omega\) extends the calculation to all their smooth combinations and proves
\begin{equation}
 \left.\mathcal F^{\rm can}_{X,Y}\right|_{\mathcal Z}=0
 \quad\hbox{for }X,Y\hbox{ characteristic}.
 \label{hmtcuatro:cierre-caracteristico}
\end{equation}
For two lapses, \([X,Y]\) includes the tangential deformation and the Gauss actions of \eqref{hmtcuatro:algebra-restricciones}. On passing to the gauge quotient, the equality is expressed modulo these vertical actions. It is not asserted that \(\Omega\) is zero on the whole phase space or that every global holonomy of a leaf is trivial.

\par\medskip\noindent\textbf{7. Prequantum expression of the same closure.}\quad
The differential realization associated with \eqref{hmtcuatro:conexion-canonica}, on smooth sections where the products are defined, is
\begin{equation}
 Q_{\rm pre}(f)=-i\hbar\nabla_{X_f}+f
            =-i\hbar X_f+f-\Theta(X_f).
 \label{hmtcuatro:operador-precuantico}
\end{equation}
With \(R^\nabla=(i/\hbar)\Omega\), expansion of the commutator and \(X_f(g)=-\{f,g\}_D\) give
\begin{equation}
 [Q_{\rm pre}(f),Q_{\rm pre}(g)]
           =i\hbar Q_{\rm pre}(\{f,g\}_D).
 \label{hmtcuatro:representacion-precuantica}
\end{equation}
The product rule preserving the structure functions is
\begin{equation}
 Q_{\rm pre}(aC)
   =aQ_{\rm pre}(C)+C\bigl(Q_{\rm pre}(a)-a\bigr).
 \label{hmtcuatro:funciones-estructura}
\end{equation}
It follows from \(X_{aC}=aX_C+CX_a\) and \eqref{hmtcuatro:operador-precuantico}. The second term is not removed outside \(\mathcal Z\). On \(\mathcal Z\), \(Q_{\rm pre}(C_A)=-i\hbar\nabla_{X_{C_A}}\); thus \eqref{hmtcuatro:representacion-precuantica} and \eqref{hmtcuatro:cierre-caracteristico} are the operator and geometric expressions of the same canonical closure, not two independent incorporations of gravity.
\end{proof}

\subsection{Domain of the result and comparison of realizations}
\label{hmtcuatro:alcance-canonico}

The proof brings together an action with dynamical gravity and a joint source, its Legendre transform and its characteristic connection. The qualifier \emph{arbitrary} for deformations refers to admissible lapses and shifts generating characteristic directions in that regular region; not to any phase-space vector, any boundary or any singular solution. Local flatness preserves possible global monodromy and does not identify two TPK histories merely because they have the same endpoints.

The coefficient \(H_X^{\rm can}\) of a line connection is not by definition the clock-and-response operator \(H_{\rm clk}+D^*WD\), or its interacting extension. \(Q_{\rm pre}(C)\) additionally includes a derivation on phase space; it is not an alternative name for those operators either. Identifying a concrete realization through a map \(I\) requires the typed equality
\[
 \nabla_N^{\rm op}I=I\nabla_N^{\rm can}
\]
on their effective domains. If \(I\) depends on geometry, its derivative is part of this equality.

The distinction has an exact check: in the chart \(\Theta=p\,\mathrm dq\),
\begin{equation}
 Q_{\rm pre}(p^2/2)=-i\hbar p\,\partial_q-p^2/2.
 \label{hmtcuatro:control-polarizacion}
\end{equation}
Acting on the function \(1\), it produces \(-p^2/2\). It therefore does not preserve by mere restriction the polarization of functions independent of \(p\), and it is not \(-\hbar^2\partial_q^2/2\) there. The check does not refute the received Hamiltonian: it prevents replacing one representation by another without its map.

No infinite-dimensional Liouville measure has been assumed, nor is self-adjointness deduced from the preceding differential identities. If refinement inclusions intertwine the connections and preserve a common core, their curvatures intertwine by composition; a limit additionally requires its own domain and convergence controls. That condition is not presented here as verified for the interacting operator and the joint spatial--chiral limit. The earlier quantum constructions are preserved; the present result specifies exactly the canonical closure just proved.
